\documentclass[a4paper]{report}
% Shared preamble for the English probability and statistics documents.
% Based on the original course preamble. Compile with pdfLaTeX.
\usepackage[top=1.5cm,bottom=2cm,left=2cm,right=2cm]{geometry}
\usepackage[T1]{fontenc}
\usepackage{lmodern}
\usepackage[utf8]{inputenc}
\usepackage[english]{babel}

\usepackage{amsmath}
\usepackage{amssymb}
\usepackage{amsthm}
\usepackage{mathrsfs}
\usepackage{mathtools}
\usepackage{etoolbox}
\usepackage{tikz}
\usetikzlibrary{shapes.multipart,arrows.meta}
\usepackage{xcolor}
\usepackage{enumerate}
\usepackage[shortlabels]{enumitem}
\usepackage{hyperref}
\usepackage{version}
\usepackage{bbold}
\usepackage{fontawesome5}
\usepackage{fancyhdr}
\lfoot{}
\usepackage{natbib}
\usepackage{graphicx}
\usepackage{pdfpages}
\usepackage{authblk}
    
\usepackage{blindtext}
\usepackage{multicol}
\usepackage{collcell}
\usepackage{wrapfig}
\usepackage{blkarray}
\usepackage{cancel}
\usepackage{appendix}
\usepackage{float}
\makeatletter
\input{tex/assets/slashbox.sty}
\makeatother
\hypersetup{unicode,colorlinks,citecolor=red,filecolor=black,
  linkcolor=red,urlcolor=cyan}
\setlength{\columnsep}{1cm}
\pagenumbering{gobble}

\def\numtext#1{%
  \ifcase#1\or one\or two\or three\or four\or five\or six\or seven\or eight\or nine\or ten\or
  eleven\or twelve\or thirteen\or fourteen\or fifteen\or sixteen\or seventeen\or eighteen\or nineteen\or twenty\or Lots\fi}

\newcommand{\Byte}[1]{%
  \def\nodecontents{}%
  \foreach[count=\l] \k in {#1} {\xappto{\nodecontents}{\noexpand\nodepart{\numtext{\l}}\k}}%
  \node[rectangle split,rectangle split parts=8,rectangle split horizontal,draw] at (2,2){\nodecontents};}

\newcommand{\D}{{\mathrm{d}}}
\newcommand{\cE}{{\mathcal{E}}}
\newcommand{\cG}{{\mathcal{G}}}
\newcommand{\cX}{{\mathcal{X}}}
\newcommand{\cY}{{\mathcal{Y}}}
\newcommand{\cB}{{\mathcal{B}}}
\newcommand{\cF}{{\mathcal{F}}}
\newcommand{\cO}{{\mathcal{O}}}
\newcommand{\cS}{{\mathcal{S}}}
\newcommand{\cT}{{\mathcal{T}}}
\newcommand{\given}{\, \big| \,}
\newcommand{\cL}{{\mathscr{L}}}
\newcommand{\cD}{{\mathcal{D}}}
\newcommand{\cA}{{\mathscr{A}}}
\newcommand{\cP}{{\mathscr{P}}}
\newcommand{\cR}{{\mathscr{R}}}
\newcommand{\A}{{\mathcal A}}
\newcommand{\one}{\mathbb{1}}
\newcommand{\mathbbm}[1]{\mathbb{#1}}

\providecommand{\C}{}
\renewcommand{\C}{\mathbb C}
\newcommand{\R}{\mathbb R}
\newcommand{\Q}{\mathbb Q}
\newcommand{\N}{\mathbb N}
\newcommand{\T}{\mathbb T}
\newcommand{\F}{\mathbb F}
\renewcommand{\P}{\mathbb P}
\newcommand{\Z}{\mathbb Z}
\newcommand{\E}{\mathbb {E}}
\newcommand{\Pbb}[1]{\P\left( #1\right)}
\newcommand{\Ebb}[1]{\E\left[ #1\right]}

\newcommand{\lb}{\left (}
\newcommand{\rb}{\right )}
\newcommand{\lbb}{\left [}
\newcommand{\rbb}{\right ]}
\newcommand{\labs}{\left |}
\newcommand{\rabs}{\right |}
\newcommand{\lbrb}[1]{\lb #1 \rb}
\newcommand{\lbbrbb}[1]{\lbb#1\rbb}
\newcommand{\labsrabs}[1]{\labs#1\rabs}
\newcommand{\lbbrb}[1]{\lbb#1\rb}
\newcommand{\lbrbb}[1]{\lb#1\rbb}
\newcommand{\langlerangle}[1]{\langle#1\rangle}
\newcommand{\lbcurly}{\{}
\newcommand{\rbcurly}{\}}
\newcommand{\curly}[1]{\lbcurly#1\rbcurly}
\newcommand{\abs}{\labsrabs}
\newcommand{\ind}[1]{1_{#1}}
\newcommand{\indep}{\perp \!\!\! \perp}
\renewcommand{\d}{\, \mathrm{d}}

\newtheorem{theorem}{Theorem}
\newtheorem{problemCnt}{Lemma}
\theoremstyle{definition}
\newtheorem{problem}[problemCnt]{Problem}
\newtheorem{problembook}[problemCnt]{\faBook \,  Problem}
\newtheorem{problempencil}[problemCnt]{\faPencil* \,  Problem}
\newtheorem{problemstar}[problemCnt]{\faStar\,  Problem}
\newtheorem{problemstarpencil}[problemCnt]{\faStar \, \faPencil* \,   Problem}
\newtheorem{probleminteresting}[problemCnt]{\faSmile[regular] \,   Problem}
\newtheorem{solution}{Solution}

\newcommand{\fwcell}[1]{\makebox[\arraystretch\normalbaselineskip]{$#1$}}
\newenvironment{enumeratei}{\begin{enumerate}[(i)]}{\end{enumerate}}
\newenvironment{enumeratea}{\begin{enumerate}[(a)]}{\end{enumerate}}
\newenvironment{enumeraten}{\begin{enumerate}[\upshape 1.]}{\end{enumerate}}

% Compatibility helpers needed by a few files in the combined archive.
\AtBeginDocument{\noindent{\footnotesize English translation prepared with ChatGPT (OpenAI, GPT-6), 26 October 2026, for translation of the Bulgarian text into English and LaTeX formatting.
The Bulgarian originals remain the reference; translation may contain errors.}\par\smallskip}
\newcommand{\MissingGraphic}[1]{%
  \begin{center}\fbox{\parbox{0.82\linewidth}{\centering Missing source image: \texttt{#1}}}\end{center}}



\newif\ifshowsolutions
\showsolutionstrue
\begin{document}
\noindent Probability and Statistics, CS Year 3, 2020/2021 \hfill Test 1\ifshowsolutions: solutions\fi
\hrulefill
\begin{it}Your grade is 2 + the number of points. Time allowed: 3 hours. Good luck! Work on $(\Omega,\mathcal F,\P)$.\end{it}
\begin{problem} (1 point)
At a shooting gallery, a boy pays 2 BGN for three shots. The prizes are 10 BGN for three hits, 5 for two hits, and 1 for one hit. His hit probability per shot is $1/3$. Is the game fair?\footnote{A game is fair if expected net winnings are zero.} For which values of the largest prize is the game favourable to him?
\ifshowsolutions
\par\textbf{Solution.} (0.6 points for the main part + 0.1 for $y>12$.) Let the largest prize be $y$. Then
\begin{align*}
\E X={}&-2\binom30(2/3)^3-\binom31(1/3)(2/3)^2\\
&+3\binom32(1/3)^2(2/3)+(y-2)\binom33(1/3)^3=\frac{y-2}{27}.
\end{align*}
Thus the game is favourable for $y>12$, unfavourable for $y<12$, and in particular is not fair when $y=10$.
\fi
If the boy makes a profit, what is the probability he hit the target on his third shot?
\ifshowsolutions
\par\textbf{Solution.} (0.3 points.) Write $Y$ for a hit and $N$ for a miss. The probability of missing the third shot is
\[\P(YYN\mid\{YYN,YNY,NYY,NNN\})=\frac{2/27}{3\cdot2/27+1/27}=\frac27.\]
The answer is therefore $1-2/7=5/7$.
\fi
\end{problem}
\begin{problem}
The first box contains 4 white and 6 black balls; the second 7 white and 3 black. Draw one ball randomly from each, add one white ball, and place all three in a third box.
\begin{enumerate}
\item (0.4 points) What is the probability that a random ball from the third box is white?
\ifshowsolutions
\par\textbf{Solution.} Let $A$ mean a white draw from the third box, and $B_{xy}$ mean colours $x,y$ from the first two, with $W$ for white and $B$ for black. Total probability gives
\begin{align*}
\P(A)&=\P(A\mid B_{WW})\P(B_{WW})+\P(A\mid B_{BW})[\P(B_{BW})+\P(B_{WB})]+\P(A\mid B_{BB})\P(B_{BB})\\
&=1\cdot\frac{4\cdot7}{100}+\frac23\frac{4\cdot3+6\cdot7}{100}+\frac13\frac{6\cdot3}{100}=\frac7{10}.
\end{align*}
\fi
\item (0.6 points) Draw from the third box with replacement until the first white ball. Let $X$ count draws. Find $\E X$. Is $X$ geometric?
\ifshowsolutions
\par\textbf{Solution.} (0.3 points if all calculations are otherwise correct but use a geometric law with the probability from part 1.) With three white balls the first draw succeeds. With two white and one black, $X\sim\mathrm{Ge}(2/3)$ conditionally and its mean is $3/2$. With one white and two black the conditional law is $\mathrm{Ge}(1/3)$, with mean 3. Thus
\begin{align*}
\E X&=\E[X\mid B_{WW}]\P(B_{WW})+\E[X\mid B_{BW}][\P(B_{BW})+\P(B_{WB})]+\E[X\mid B_{BB}]\P(B_{BB})\\
&=\frac{28}{100}+\frac32\frac{54}{100}+3\frac{18}{100}=\frac{163}{100}.
\end{align*}
If $X$ were geometric, part 1 would require $\mathrm{Ge}(7/10)$ and mean $10/7$, contrary to this calculation.
\fi
\end{enumerate}
\end{problem}
\begin{problem} (1 point)
A bag has $n\in\mathbb N$ coins with values $1,\ldots,n$. A person takes a handful, with each coin selected with probability $p\in(0,1)$. Find the expected sum $S$.
\ifshowsolutions
\par\textbf{Solution.} (0.1--0.2 or more for partial reasoning.) Let $X_k$ indicate selection of coin $k$. Then $X_k\sim\mathrm{Ber}(p)$ and $\E X_k=p$. Linearity gives
\[\E S=\E\left[\sum_{k=1}^nkX_k\right]=\sum_{k=1}^nk\E X_k=\frac{pn(n+1)}2.\]
\fi
\end{problem}
\begin{problem}
\begin{enumerate}
\item (0.4 points) State when events $A,B\subset\Omega$ are independent and when they are disjoint. Can they be independent if (a) disjoint, or (b) $A\subset B$? If so, under what conditions?
\ifshowsolutions
\par\textbf{Solution.} Award 0.1 for each definition and each example. Disjoint means $A\cap B=\emptyset$; independent means $\P(A\cap B)=\P(A)\P(B)$. Both answers are yes: (a) requires $\P(A)=0$ or $\P(B)=0$; (b) requires $\P(A)=0$ or $\P(B)=1$.
\fi
\item (0.6 points) State when discrete random variables $X,Y$ are independent and when identically distributed.
\ifshowsolutions
\par\textbf{Marking.} 0.1 per definition, 0.1 for $\P(X<Y)=1/2$, 0.1 for inclusion--exclusion, with at most 0.3 for these first four items; 0.3 for the last part.
\par\textbf{Solution.} Identical distributions mean $\P(X=a)=\P(Y=a)$ for every $a\in\mathbb R$. Independence means $\P(X=a,Y=b)=\P(X=a)\P(Y=b)$ for all $a,b$.
\fi
Let $X,Y,Z$ be identically distributed and pairwise independent. Assume ties have probability zero. Find $\P(X<Y)$, $\P(Y<Z)$ and $\P(Z<X)$.
\ifshowsolutions
\par\textbf{Solution.} By independence and identical laws, $\P(X<Y)=\P(X>Y)$. Since ties have probability zero and the three comparison probabilities sum to one, each requested probability is $1/2$.
\fi
State inclusion--exclusion for three events.
\ifshowsolutions
\begin{align*}
\P(A\cup B\cup C)={}&\P(A)+\P(B)+\P(C)\\
&-\P(A\cap B)-\P(A\cap C)-\P(B\cap C)+\P(A\cap B\cap C).
\end{align*}
\fi
Prove $a=\min\{\P(X<Y),\P(Y<Z),\P(Z<X)\}\leq2/3$.\footnote{Can you construct equality? A simple example uses suitable probabilities on events of the form $(X=n_1,Y=n_2,Z=n_3)$.}
\ifshowsolutions
\par\textbf{Solution.} Set $A=\{X<Y\}$, $B=\{Y<Z\}$, $C=\{Z<X\}$. Their triple intersection is empty. Their pairwise intersections are disjoint, so
\[\P(A\cap B)+\P(A\cap C)+\P(B\cap C)=\P((A\cap B)\cup(A\cap C)\cup(B\cap C))\leq1.\]
Inclusion--exclusion gives $1\geq3a-1$, as required.
\par\emph{This illustrates how, whichever of three candidates wins, two thirds of voters may prefer another candidate. Three voters can have rankings $A>B>C$, $B>C>A$, $C>A>B$. In probabilistic notation,
$\P(X=1,Y=2,Z=3)=\P(X=2,Y=3,Z=1)=\P(X=3,Y=1,Z=2)=1/3$.
A related observation dates to Condorcet in the late eighteenth century; see the Condorcet paradox.}
\fi
\end{enumerate}
\end{problem}
\end{document}
